\section{The actual host character frame and the original column input}
\label{sec:host-character-frame}

\status{HF1--HF13 construct the five actual host character families,
with their prescribed infinity types, restrictions, tame local
characters and inertia orders. HF14--HF15 construct their determinant
quotient and relative determinant laws. HF16--HF19 join these
constructions to the previously reconstructed original geometry,
coefficient fields, column laws, jets and later stages on the same
base. The final paragraph identifies precisely why these existence
results do not yet construct a quantitative \texttt{V4.Core}.}

We use the literal host of the source. Write $E$ for its CM field,
$F=E^+$ for its maximal real subfield, $K=\mathbb Q(\sqrt{-d})$
for its distinguished imaginary quadratic subfield, and
$c:x\mapsto\bar x$ for its specified conjugation. The CM-type is
$\Phi=(\phi_\tau)_{\tau\in\Psi}$, where
$\phi_\tau(\sqrt{-d})=i\sqrt d$. AH9--AH12 establish these
identifications and the fact that $E/F$ is unramified at every
finite prime. AC2--AC3 establish the host's root-of-unity
identities. A general CM field used in intermediate arguments
will be denoted by $L$.

Let $\mathbb I_L$ be the full idele group and
\[
 U_L=L_\infty^\times\times\prod_P\mathcal O_{L_P}^{\times}
       \subset\mathbb I_L,\qquad
 \operatorname{ph}(z)=z/|z|\quad(z\in\mathbb C^\times).
 \tag{HF.1}
\]
This $U_L$ is the source's \texttt{unitIdl}; its infinite
components are unrestricted. A Hecke character is a continuous
homomorphism $\mathbb I_L/L^\times\to\mathbb C^\times$.
Unitary means that every value has absolute value one.
Our finite valuations $\operatorname{ord}_P$ are additive with a
uniformizer of order one; where the source uses multiplicative
valuations, a principal unit satisfies $|u-1|_P<1$.
The exact infinity convention is
\[
 \chi_\infty(k)=\prod_\tau
       \operatorname{ph}(\phi_\tau(k))^{n_\tau}.
 \quad
 \texttt{HasInfQ}(a)\ \text{means}\ n_\tau=2a_\tau\in\mathbb Z.
 \tag{HF.2}
\]
In particular the factor two in the definition of
\texttt{IsMu0.inf} is retained throughout.

\paragraph{HF1: roots of unity and the compatibility needed for extension.}
For a CM field $L/F$, every embedding carries $c$ to ordinary
complex conjugation. If $x\bar x=1$ and
$\operatorname{ord}_P(x)=0$ at every finite prime, then $x$ is
an algebraic integer and every complex conjugate of $x$ has
absolute value one. To see directly that it is a root of unity,
consider the monic characteristic polynomials over $\mathbb Q$
of $x,x^2,x^3,\ldots$ acting on $L$. Their coefficients are
integers, and their elementary-symmetric formulas give uniform
bounds depending only on $[L:\mathbb Q]$. There are only finitely
many such polynomials and finitely many roots of all of them.
Two positive powers of $x$ coincide, so $x$ has finite order.
Conversely a root of unity is an algebraic-integer unit.

For the actual host, AC2--AC3 show that the specified element
$s_g\in\Psi$ is a nonidentity involution and
$s_g(\zeta)=\zeta^{-1}$ for every root of unity $\zeta\in E$.
The map $\tau\mapsto\tau s_g$ pairs the CM-type without fixed
points. The host group is abelian, so
\[
 \phi_{\tau s_g}(\zeta)=\phi_\tau(\zeta)^{-1},\qquad
 \prod_{\tau\in\Psi}\phi_\tau(\zeta)=1,\qquad
 N_{E/K}(\zeta)=1.                                    \tag{HF.3}
\]
The final equality also follows before applying the distinguished
complex embedding: $\Psi$ is exactly $\operatorname{Gal}(E/K)$,
and the same pairing cancels its norm product.

For an integral weight $e:\Psi\to\mathbb Z$, the source's
$\mathrm{MU}(\Phi,e)$ is the assertion
$\prod_\tau\phi_\tau(\zeta)^{e_\tau}=1$ for every such $\zeta$.
The predicate $\mathrm{MUQ}$ includes a choice of these integral
weights equal to the specified rational weights. Thus it gives
the same assertion with exponent $-e$ by inversion. We will use
the actual $\mathrm{MUQ}$ hypothesis for nonconstant weights;
(HF.3) proves it for constant integral weights.

\paragraph{HF2: continuous circle characters with prescribed values.}
Let $G$ be an abelian topological group, $U\subset G$ an open
subgroup, $\psi:U\to S^1$ continuous, $\alpha:A\to G$ a
homomorphism from an abelian group, and $a:A\to S^1$ another
homomorphism. Suppose
\[
       \alpha(y)\in U\quad\Longrightarrow\quad
                    a(y)=\psi(\alpha(y)).             \tag{HF.4}
\]
On $U\alpha(A)$ set
$\gamma(u\alpha(y))=\psi(u)a(y)$.
If two expressions represent the same element, their quotient
is an element of $\alpha(A)\cap U$, and (HF.4) proves equality
of their values. This is a homomorphism, continuous because its
restriction to the open subgroup $U$ is $\psi$.

Every homomorphism from a subgroup of an abelian group to $S^1$
extends to the entire group. Here is the algebraic argument
used in this application. Order partial extensions by inclusion
and take a maximal one. For an element $g$ outside its domain,
if no positive power of $g$ lies in the domain, choose any value
of absolute value one and extend on the subgroup generated by
$g$. Otherwise the integers $n$ with $g^n$ in the domain form
$n_0\mathbb Z$. Choose an $n_0$th root in $S^1$ of the known
value at $g^{n_0}$, and define the extension on $h g^j$
multiplicatively. The minimality of $n_0$ makes it well defined.
Both cases contradict maximality unless the domain was all of
$G$. The resulting extension $\nu:G\to S^1$ is continuous:
continuity at one is already known on the open subgroup $U$,
and translations give continuity everywhere.
Consequently
\[
 \nu|_U=\psi,\qquad \nu(\alpha(y))=a(y).
 \tag{HF.5}
\]
This proves the substantive open-subgroup extension used by
both character producers. It uses divisibility of the circle,
not a global reciprocity existence theorem.

\paragraph{HF3: the finite-unit model and its exact local readings.}
For a finite prime $P$ of $L$, choose a uniformizer $\pi_P$ in
$L_P$, and for $x\in L_P^\times$ put
\[
       x^\circ=x\pi_P^{-\operatorname{ord}_P(x)}.
 \tag{HF.6}
\]
This is a multiplicative map into the units of the local
valuation ring. If $\theta_P:\kappa(P)^\times\to S^1$ is a
residue character, then
$\widetilde\theta_P(x)=\theta_P(\overline{x^\circ})$
is a continuous unitary character of $L_P^\times$.
It kills $1+P\mathcal O_{L_P}$ and agrees with $\theta_P$
on units. A residue of a local unit may be computed from any
global algebraic integer within distance strictly less than
one; two such approximations differ by an element of $P$.
Thus this description agrees with the source's
\texttt{resPart}, including its readings on rational integers.
Every nonzero residue class is reached by an algebraic integer
which is a unit at $P$, so these readings also detect the exact
order of $\theta_P$.

For a CM-type $\phi_i$ selecting each infinite place once,
integers $n_i$, and a finite set $S$ of finite primes, define
\[
 \psi(z)=\prod_i\operatorname{ph}(z_{\phi_i})^{n_i}
                \prod_{P\in S}\widetilde\theta_P(z_P).
 \tag{HF.7}
\]
The completion at $\phi_i$ is identified with $\mathbb C$ by
that specified embedding. Finite products and the local
continuity just proved make $\psi$ continuous and unitary.
At an infinite component it has exactly the stated phase;
at a finite component it is $\widetilde\theta_P$ if $P\in S$,
and is one otherwise. In particular
\[
 \psi(\operatorname{diag}k)
   =\prod_i\operatorname{ph}(\phi_i(k))^{n_i}
      \prod_{P\in S}\widetilde\theta_P(k).
 \tag{HF.8}
\]
The subgroup $U_L$ is open: it is the product of the infinite
group with the standard open compact unit subgroup of the
finite ideles. Imposing finitely many principal-unit
congruences remains open.

Conjugation extends isometrically from $L$ to
$L_P\to L_{\bar P}$ and to all archimedean completions.
It therefore defines a continuous involution $C$ of the
idele group, preserves $U_L$, takes a local component at $P$
to the conjugate component at $\bar P$, and fixes every
idele coming from $F$. These assertions follow coordinate by
coordinate; no identification of different selected primes is
made. They are the actual \texttt{IsConjI} laws used below.

\paragraph{HF4: anticyclotomic characters with a prescribed tame part.}
Let $L/F$ be CM with CM-type $\phi_i$. Let $e_i\in\mathbb Z$,
let $S$ contain no pair $P,\bar P$, including no fixed prime,
and prescribe residue characters $\theta_P$ at $P\in S$.
Assume the exact root-of-unity compatibility
\[
 \prod_i\phi_i(\zeta)^{-e_i}
       \prod_{P\in S}\theta_P(\bar\zeta_P)=1
                 \quad\text{for every }\zeta\in\mu(L).
 \tag{HF.9}
\]
Use (HF.7) with $n_i=-e_i$. Let $A$ be the subgroup of
$L^\times$ consisting of $x$ with $x\bar x=1$, map it
diagonally to the ideles, and prescribe value one on $A$.
If its diagonal lies in $U_L$, HF1 makes $x$ a root of
unity, so (HF.9) is exactly the compatibility (HF.4).
HF2 gives a continuous unitary $\nu$ agreeing with $\psi$
on $U_L$ and killing these norm-one principal ideles. Set
\[
                    \chi(z)=\nu(z/Cz).               \tag{HF.10}
\]
For $k\in L^\times$, $k/\bar k$ has norm one. Thus
$\chi(\operatorname{diag}k)=1$, and $\chi$ is a Hecke
character. It is unitary and continuous. It is trivial
on the ideles of $F$ because $C$ fixes them.

For an infinite component,
$\operatorname{ph}(z/\bar z)=\operatorname{ph}(z)^2$;
therefore its infinity exponents are exactly $-2e_i$.
For a finite local unit $u$ at $P$, the idele
$u/Cu$ lies in $U_L$, so it can be evaluated using $\psi$.
The absence of conjugate pairs in $S$ gives
\[
 \begin{array}{ll}
  \chi_P(u)=\theta_P(\bar u),
                         &P\in S,\\
  \chi_{\bar P}(u)=
       \theta_P(\overline{c(u)})^{-1},
                         &P\in S,\\
  \chi_Q(u)=1,             &Q\notin S\cup\bar S.
 \end{array}                                                \tag{HF.11}
\]
Every finite local character is consequently tame.
For every nonnegative integer $n$,
\[
 \bigl[\chi_P(u)^n=1\text{ for every local unit }u\bigr]
   \Longleftrightarrow \theta_P^n=1
   \Longleftrightarrow
 \bigl[\chi_{\bar P}(u)^n=1\text{ for every local unit }u\bigr].
 \tag{HF.12}
\]
Surjectivity of reduction proves the forward implications;
conjugation is a local-field isomorphism, and inversion does
not change the kernel of a power. This is the exact local-order
part of \texttt{exists\_hecke\_trivOn}, as well as its
unitarity, real-subfield restriction and infinity-type clauses.

\paragraph{HF5: finite order when the infinity type is trivial.}
Suppose a continuous Hecke character $\chi$ is tame at a
finite set $S$, unramified away from $S$, and has trivial
infinity type. Its infinite restriction is identically one:
the diagonal $L$ is dense in $L\otimes_{\mathbb Q}\mathbb R$
by a rational basis and density of $\mathbb Q$ in $\mathbb R$,
so the prescribed identity on nonzero diagonal elements
extends by continuity to the archimedean unit group.

Choose a positive integer $M$ killing all the residue
characters at $S$, for example the product of
$|\kappa(P)^\times|$. Then $\chi^M$ kills every finite
local unit and the full infinite group. Finite-support unit
ideles are dense in the product of finite local unit groups;
hence continuity shows that $\chi^M$ kills all of $U_L$.
The finite valuations give
\[
       \mathbb I_L/(U_L L^\times)\simeq
                                  \operatorname{Cl}(\mathcal O_L).
 \tag{HF.13}
\]
Indeed their map onto fractional ideals is onto by products
of local uniformizers, its kernel is $U_L$, and diagonal
elements give precisely principal ideals. The class group is
finite by the lattice proof in TL14. If its order is $h$, then
$\chi^{Mh}=1$. If every local unit character already has order
dividing $15$, we may take $M=15$.
This proves the finite-order supplier required below; exact
order $15$ on inertia does not assert exact global order $15$.

\paragraph{HF6: selected split primes and the residue norm product.}
If a rational prime $p$ has $[E:\mathbb Q]$ distinct primes
above it, the identity $\sum_{P\mid p}e_Pf_P=[E:\mathbb Q]$
forces every $e_P=f_P=1$. Thus
$\kappa(P)=\mathbb F_p$ by the unique prime-field map,
$p\notin P^2$, and the decomposition group of every such
prime is trivial. Conjugation is nontrivial, so
$P\ne\bar P$. The two primes of $K$ above $p$ are exchanged
by conjugation.

Choose one prime $\mathfrak p$ of $K$ above $p$, and let
$S(\mathfrak p)$ be all primes of $E$ above it.
The group $\operatorname{Gal}(E/K)$ acts simply transitively
on this set. Fix $P_0\in S(\mathfrak p)$ and, for every $P$,
choose the unique automorphism carrying $P$ to $P_0$.
Since all residue maps on rational integers are the same,
reducing the norm product at $P_0$ gives
\[
       \prod_{P\in S(\mathfrak p)}\bar u_P
                    =\overline{N_{E/K}(u)}
                       \quad\text{in }\mathbb F_p^\times
 \tag{HF.14}
\]
for every $u\in\mathcal O_E^\times$. Therefore, if a fixed
character $\omega:\mathbb F_p^\times\to S^1$ is transported
to each of these residue fields and $N_{E/K}(u)=1$,
the product of the resulting character values is one.
This proves the exact residue-product supplier used for both
$\mu_0$ and $\delta$, including its uniformity over every
prime above the selected quadratic prime.

For each of the two auxiliary primes $g_k$, lying-over chooses
such a prime $\mathfrak p_k$ of $K$. The finite multiplicative
group $\mathbb F_{g_k}^\times$ is cyclic: a finite subgroup of
a field has at most $a$ roots of $X^a-1$, which forces its
exponent to equal its order. Choose a generator and send it
to a primitive complex $(g_k-1)$st root. This constructs
the injective character $\omega_k$ and hence the literal
$P:\texttt{PiBar}$ and $W:\texttt{Omg}$.
More generally, if $r\mid p-1$, sending a generator to a
primitive $r$th root constructs a character $\omega$ with
$\omega^n=1$ exactly when $r\mid n$.

\paragraph{HF7: the actual $\mu_{0,j}$ characters.}
Fix the chosen $P,W$ and an integral weight
$e_j:\Psi\to\mathbb Z$ satisfying $\mathrm{MU}(\Phi,e_j)$.
The two distinct auxiliary primes $g_0,g_1$ have disjoint
sets of primes above them. Put
$S=\bigcup_k S(\mathfrak p_k)$ and, on each selected prime
above $g_k$, prescribe
\[
             \theta_P=\omega_k^{-2(j+1)}.             \tag{HF.15}
\]
No conjugate pair lies in $S$. For a root of unity $\zeta$,
the infinite part of (HF.9) is one by $\mathrm{MU}$, and
each finite product is one by (HF.3) and (HF.14).
HF4 therefore constructs a character $\mu_{0,j}$ which
is unitary, is trivial on $F$, has infinity exponents
$-2e_{j,\tau}$, is tame everywhere, and is unramified
away from $g_0,g_1$. On every $P\mid\mathfrak p_k$ and
every natural integer $a$ prime to $g_k$ it satisfies
\[
       \mu_{0,j,P}(a)
             =\omega_k(a\bmod g_k)^{-2(j+1)}.          \tag{HF.16}
\]
These are all six fields of the literal
$\texttt{IsMu0}(P,W,e_j,j)$.
In particular its \texttt{HasInfQ} weight is $-e_j$,
not $-2e_j$. If the weights were supplied rationally,
$\mathrm{MUQ}$ supplies their integral representatives.
Doing this for every $j\in\{0,\ldots,m-1\}$ uses the same
$P,W$ throughout.

\paragraph{HF8: the unramified CM-type character $\chi_\Psi$.}
In HF4 take $S=\varnothing$ and $e_\tau=-1$ for every
$\tau$. Condition (HF.9) is exactly the middle equality
of (HF.3). The resulting $\chi_\Psi$ is unitary, trivial
on the ideles of $F$, unramified at every finite prime,
and has infinity exponent $2$ at every member of the
actual CM-type. This proves all four fields of
$\texttt{IsChiPsi}(H,\chi_\Psi)$.
The construction is on the actual host; it does not replace
the host by an imaginary quadratic field or by the synthetic
character types of a bare local-stage model.

\paragraph{HF9: the actual quadratic product formula on the real subfield.}
For a finite prime $v$ of $F$, put
\[
 \eta_v(a)=
 \begin{cases}
     (-1)^{\operatorname{ord}_v(a)},&v\text{ inert in }E/F,\\
     1,&v\text{ split in }E/F,
 \end{cases}
 \qquad
 \eta_w(a)=\operatorname{sgn}(w(a))\quad(w\mid\infty).
 \tag{HF.17}
\]
There is no ramified finite case by AH11. We prove, for this
actual host, the product formula
\[
             \prod_v\eta_v(a)\prod_{w\mid\infty}\eta_w(a)=1
                         \qquad(a\in F^\times).        \tag{HF.18}
\]
All finite products here have finite support.

First, at a prime $v\nmid d$ of residue cardinality $p^f$,
the sign distinguishing inert and split is
\[
             \eta_v(\pi_v)=\left(\frac{p^f}{d}\right).
 \tag{HF.19}
\]
For odd $p$, reduce the integral element $s_d=\sqrt{-d}$
at a prime $P$ over $v$. It is nonzero. The arithmetic
Frobenius of the unramified quadratic extension acts as
identity in the split case and as conjugation in the inert
case. On this element it acts by
$s_d^{p^f}=((-d/p))^f s_d$, by Euler's criterion and
iteration. Since $p\ne2$, the two signs are distinct.
Quadratic reciprocity and $d\equiv3\pmod4$ identify
$(-d/p)$ with $(p/d)$. At $p=2$, AH12 gives splitting,
and $d\equiv7\pmod{16}$ gives $(2/d)=1$, so (HF.19)
also holds there.

Let $y\in\mathcal O_F$ be nonzero and a unit at every
prime over $d$. Factoring $(y)$ and using (HF.19) gives
\[
 \prod_v\eta_v(y)=
       \left(\frac{|N_{F/\mathbb Q}(y)|}{d}\right),\qquad
 \prod_{w\mid\infty}\eta_w(y)
                    =\operatorname{sgn}N_{F/\mathbb Q}(y).
 \tag{HF.20}
\]
As $(-1/d)=-1$, their product is
$(N_{F/\mathbb Q}(y)/d)$.
We now show that this Jacobi symbol is one using an actual
quadratic subfield of $F$, rather than assuming the product
formula. The specified elements $\sqrt{-q}$ and $\sqrt{-d}$
give $\sqrt{qd}\in F$. The integer $qd$ is not a rational
square since $q,d$ are distinct primes. Put
$F_0=\mathbb Q(\sqrt{qd})$ and
$z=N_{F/F_0}(y)$. This is integral, so its trace $A$ and norm
$N=N_{F/\mathbb Q}(y)$ are integers. Writing
$2z=A+B\sqrt{qd}$ with $B\in\mathbb Q$ yields
\[
                       A^2-4N=qdB^2.                 \tag{HF.21}
\]
The right side is an integer. Its $d$-adic order is
$1+2\operatorname{ord}_d(B)$, so integrality forces
$\operatorname{ord}_d(B)\ge0$ and hence $d\mid A^2-4N$.
The norm $N$ is prime to $d$: an inverse to $y$ modulo
$d\mathcal O_F$ exists, and determinants of multiplication
show that $N$ is a unit modulo $d$. Thus
$N\equiv(A/2)^2\pmod d$ is a nonzero square, proving
(HF.18) for $y$ prime to $d$.

To remove this restriction, let $y\in\mathcal O_F$ be arbitrary
and nonzero. Every prime $v\mid d$ splits in $E/F$ by AH12.
For each pair $P,\bar P$ above such a $v$, use CRT modulo
sufficiently large powers of these distinct primes to choose
$b\in\mathcal O_E$ with
\[
 \operatorname{ord}_P(b)=\operatorname{ord}_v(y),
                    \qquad\operatorname{ord}_{\bar P}(b)=0.
 \tag{HF.22}
\]
The exact order $n$ is imposed by congruence to a local
uniformizer power modulo $P^{n+1}$. Let $a=b\bar b\in F$.
Unramifiedness gives
$\operatorname{ord}_v(a)=\operatorname{ord}_v(y)$ at the
chosen bad primes. Moreover the product in (HF.18) is one
for $a$: all real embeddings are positive, and at an inert
prime its order is twice the order of $b$.

Consequently $x=y/a$ is a unit at every prime above $d$.
Its denominator ideal is coprime to their product. CRT
therefore supplies $t\in\mathcal O_F$, a unit at those
primes, with $tx\in\mathcal O_F$; both $t$ and $tx$ are
units there. The already proved case applies to each of
them, hence to $x$, and then to $y=ax$. Finally write an
arbitrary nonzero element of $F$ as a quotient of two
algebraic integers. Multiplicativity proves (HF.18) in
full. This is the substantive host-specific proof of
\texttt{quadraticCharProductFormulaCMHost}.

\paragraph{HF10: an unramified character with the quadratic restriction.}
We next prove the extension principle actually used for
$\eta$. Let $L/F$ be a CM quadratic extension unramified at
finite primes, with CM-type $\phi_i$ and odd integers $n_i$.
Assume (HF.18) for its finite and infinite quadratic signs.
Suppose that for every root of unity $\zeta\in L$ one can
choose $\delta\in L^\times$ satisfying
\[
 \delta=\zeta\bar\delta,\qquad
 \operatorname{ord}_P(\delta)\in2\mathbb Z
       \quad\text{when }P\text{ lies over an inert prime of }F,
 \qquad
 \prod_i\operatorname{ph}(\phi_i(\delta))^{n_i}=1.
 \tag{HF.23}
\]
The character to be extended on $U_L$ is
$\psi(z)=\prod_i\operatorname{ph}(z_{\phi_i})^{n_i}$,
with no finite factors.

For a finite-support vector $m=(m_v)_v$ on the primes of
$F$, let $\beta(m)$ be the product of the local $F$
uniformizer ideles of exponent $m_v$, mapped to all
primes of $L$ above $v$. Set
$b(m)=\prod_v\eta_v(\pi_v)^{m_v}$.
We prescribe
\[
 \chi(\operatorname{diag}k\,\beta(m))=b(m),
          \qquad k\in L^\times.                      \tag{HF.24}
\]
It suffices, by HF2, to verify compatibility when this
idele lies in $U_L$. For $P\mid v$, its order gives
\[
                    \operatorname{ord}_P(k)=-m_v,
 \tag{HF.25}
\]
because the relative ramification index is one. Thus
$k/\bar k$ is a norm-one algebraic unit, hence a root
of unity by HF1. Choose $\delta$ for this root using
(HF.23). Then $k/\delta$ is fixed by conjugation,
so equals some $a\in F^\times$.

At an inert prime, (HF.23) and (HF.25) give
$m_v\equiv-\operatorname{ord}_v(a)\pmod2$.
The split factors are one. Hence (HF.18) implies
\[
 b(m)=\prod_v\eta_v(a)^{-1}
                =\prod_{w\mid\infty}\eta_w(a)
                =\prod_i\operatorname{ph}(\phi_i(a))^{n_i}.
 \tag{HF.26}
\]
The last equality uses oddness of every $n_i$ and the
bijection from the CM-type to the real places of $F$.
Multiplying by the last equality in (HF.23) shows
$b(m)=\psi(\operatorname{diag}k\,\beta(m))$.
This is exactly the required compatibility; HF2 now
constructs the continuous unitary character (HF.24).
Taking $m=0$ shows it kills all principal ideles.

It is unramified at every finite prime because it agrees
with $\psi$ on $U_L$. Its infinity type is $n_i$.
For a local $F$-component, divide by the corresponding
uniformizer idele of exponent $\operatorname{ord}_v(a)$;
the quotient is a finite unit idele, on which the value is
one. Equation (HF.24) then gives the finite restriction
$\eta_v(a)$. At infinity the odd phase gives
$\eta_w(a)$. Thus all clauses of the literal
\texttt{ResEta} are obtained. This proves the actual
unramified extension theorem from (HF.23), not merely
the existence of an abstract character on a quotient.

\paragraph{HF11: the host phase lemma and the actual $\eta$ character.}
We verify (HF.23) on the host with $n_\tau=1$.
Let $N=Q_Sqd$ be its specified cyclotomic conductor.
AC2--AC3 show that every host root of unity has order
dividing $N$. If $\zeta\ne-1$, set $\delta=1+\zeta$.
Then $\delta=\zeta\bar\delta$.
The involution $\tau\mapsto\tau s_g$ of HF1 takes
$\phi_\tau(\delta)$ to its complex conjugate, so their
normalized phases multiply to one. It follows that
\[
                 \prod_\tau\operatorname{ph}
                         (\phi_\tau(1+\zeta))=1.       \tag{HF.27}
\]
If $P\nmid N$ and $1+\zeta$ had positive order at $P$,
then $-\zeta$ would reduce to one. Its order divides
$N$, which is prime to the residue characteristic.
Reduction is injective on roots of such order: if
$u\ne1$, $u^a=1$ and $u\equiv1\pmod P$, the identity
$1+u+\cdots+u^{a-1}=0$ reduces to $a=0$.
This is impossible. Thus $\delta$ is a unit at every
$P\nmid N$.

For $\zeta=-1$, take $\delta=s_d=\sqrt{-d}$.
It satisfies $\delta=-\bar\delta$ and is a unit outside
$d$. At every CM-type embedding its phase is $i$.
The host rank hypotheses imply $4\mid|\Psi|$, so its
total phase is $i^{|\Psi|}=1$.
Every prime dividing $N$ splits in $E/F$ by AH12.
Consequently any prime over an inert prime of $F$
is outside $N$, and both choices of $\delta$ have
order zero there. This verifies every clause of
(HF.23), with no parity exception at residue characteristic
two.

HF9 supplies the product formula, and HF10 therefore
constructs a unitary Hecke character $\eta$, unramified
at all finite primes, with restriction $\eta_{E/F}$
on $F$ and infinity exponent one at every member of
$\Phi$. These are all four fields of
$\texttt{IsEtaT}(H,\eta)$.
This $\eta$ has the actual complex Hecke-character
domain and restriction. It is not identified with the
formal two-element character type used to inhabit a
bare local-stage record.

\paragraph{HF12: characters of exact inertial order fifteen.}
Let $p$ be a rational prime with $p\equiv1\pmod{15}$
which splits completely in $E$. Choose a prime
$\mathfrak p$ of $K$ above $p$. By HF6 there is a
character $\omega:\mathbb F_p^\times\to S^1$
of exact order $15$. Transport it to every residue
field at $P\in S(\mathfrak p)$.
Use HF4 with $e_\tau=0$ and these residue characters.
Its root compatibility follows from (HF.3) and
(HF.14). We obtain $\delta_p$, trivial on $F$,
tame at every finite prime, unramified away from $p$,
and with trivial infinity type.

Every $P\mid p$ lies either over $\mathfrak p$ or over
its conjugate. Thus (HF.12) says, literally for each
such $P$ and each $n\in\mathbb N$,
\[
 \left[\delta_{p,P}(u)^n=1
                 \text{ for every }u\in\mathcal O_{E_P}^\times\right]
                      \Longleftrightarrow 15\mid n.   \tag{HF.28}
\]
At all other primes it is unramified. HF5 gives a
positive finite global order, in fact an order dividing
$15\,|\operatorname{Cl}(\mathcal O_E)|$.
These are exactly all five fields of
$\texttt{IsDelta}(p,\delta_p)$.
Equation (HF.28) also proves $\delta_p\ne1$: take
$n=1$ at any prime above $p$.

Apply this construction at the host's distinct
$p_\delta,p_\delta'$ to obtain $\delta,\delta'$.
Their ramification supports are disjoint and nonempty,
so they are distinct from one and from each other.
If a set of forbidden rational primes omits $p$,
the same character satisfies \texttt{IsTwist} for that
set: complete splitting gives the split-prime,
prime-residue-norm and $p\notin P^2$ clauses at every
ramified prime. No \texttt{IsPsi} prime-power-order
condition is inferred; that is a different constructor.

\paragraph{HF13: the full frame on the unchanged host and weights.}
Let $H$ be the supplied host and
$e:\operatorname{Fin}(m)\to\Psi\to\mathbb Q$.
Assume that both $g_k$ and both $p_\delta,p_\delta'$
split completely in $E$, that each $e_{j,\tau}$ is
integral, and that $\mathrm{MUQ}(\Phi,e_j)$ holds for
each $j$. HF6 chooses one $P,W$; HF7 constructs all
$\mu_{0,j}$ for these same choices; HF8 and HF11
construct $\chi_\Psi$ and $\eta$; and HF12 constructs
$\delta,\delta'$. The resulting data
\[
          (P,W,\mu_0,\eta,\chi_\Psi,\delta,\delta')
 \tag{HF.29}
\]
with the proved certificates give the literal
$\texttt{HFrame}(H,e)$ of
\texttt{NonemptyHostHeckeCharacterFrame.nonempty\_hFrame}.
Each producer was proved above, including its local
and archimedean restrictions.

For a preliminary datum $P_0$ and its supplied
$\texttt{PreX}$, the required weights are
$e_{j,\tau}=(P_0.\mathrm{hc}_\tau).e_j$;
its \texttt{e\_even} is exactly the required
$\mathrm{MUQ}$ statement. Its level primes are
congruent to one modulo $N$. AH9 implies that their
inertia is trivial and their arithmetic Frobenius is
the identity; hence their decomposition groups are
trivial and they split completely. Thus HF13 gives
the actual $\texttt{Top.Fr}(P_0)$ on that same
preliminary datum. The finite field of character
values will be chosen after this frame, and orientation
of the auxiliary prime factors will be made using
this frame's $P$.

\paragraph{HF14: the actual determinant quotient and rational unitary map.}
Let $K_D$ be the subgroup generated by the following
ideles: those in $U_E$ congruent to one at every prime
above $g_0,g_1$, and the finite local ideles of elements
of $F^\times$. This is exactly the algebraic
\texttt{Subgroup.closure} in the source's \texttt{KD}.
The first set is already an open subgroup, so $K_D$
is open and closed; it consequently also equals its
topological closure. These local generators include
an element of $F^\times$ of order one at each specified
finite prime, which supplies the uniformizers used
below. The subgroup contains every infinite idele.
It also contains every diagonal
$a\in F^\times$. To prove the latter, take all real-subfield
primes below the finite set where $a$ has nonzero order or
where an auxiliary congruence is imposed. Multiply the
corresponding local $F$-ideles with component $a$.
The quotient of $\operatorname{diag}a$ by this finite product
is a unit idele equal to one at the congruence primes.
Both factors therefore lie in $K_D$.

For any character trivial on $F$, tame everywhere and
unramified away from $g_0,g_1$, its finite-idele
restriction kills $K_D$. On the first generating subgroup
this follows from tameness at the auxiliary primes and
unramifiedness elsewhere. To pass from finite products
of unit components to every unit idele, use their density
and continuity. It kills the local $F$ generators by
the exact restriction law. Its kernel is closed, so
it kills the closure. Thus the actual characters
$\mu_{0,j}$ and $\chi_\Psi$ descend to
\[
              D_T=\mathbb I_E/K_D,\qquad
              U_O=\operatorname{image}(U_E\to D_T).
 \tag{HF.30}
\]
Since $\chi_\Psi$ is unramified everywhere, its descended
finite restriction $\chi_D$ kills $U_O$.

For $z\in E^\times$ with $z\bar z=1$, an explicit
Hilbert-90 solution is
\[
 k=1+z\quad(z\ne-1),\qquad
 k=a-\bar a\quad(z=-1,\ a\ne\bar a).
 \tag{HF.31}
\]
In both cases $k\ne0$ and $z\bar k=k$.
Any two such $k$ differ multiplicatively by an element
of $F^\times$, whose diagonal lies in $K_D$.
Consequently
$e_1(z)=[\operatorname{diag}k]\in D_T$
is well defined. Products of solutions solve the
product equation, so it is a homomorphism. It has
precisely the source's rational-unitary determinant
mapping property.

The Hecke product identity on $k$ and (HF.2) for
$\chi_\Psi$ give
\[
 \chi_D(e_1(z))
   =\chi_{\Psi,\infty}(k)^{-1}
   =\prod_\tau\operatorname{ph}(\phi_\tau(k))^{-2}
   =\prod_\tau\phi_\tau(z)^{-1}.
 \tag{HF.32}
\]
For the last equality, apply $\phi_\tau$ to
$z\bar k=k$: the quotient $k/\bar k$ is the square
of its normalized complex phase. Equations
(HF.30)--(HF.32) prove the quotient, kernel,
rational-point and character laws used in the
actual determinant package.

\paragraph{HF15: relative lattice determinants and the full determinant package.}
An idele is in $K_DU_E$ if and only if its additive
orders at $P$ and $\bar P$ are equal for every conjugate
pair. One direction holds for each generator and then
for its closure, since the local order-equality conditions
are closed. Conversely, a balanced idele has only
finitely many nonzero orders. At each underlying prime
of $F$, subtract that common order using a local
$F$-uniformizer idele. Relative unramifiedness makes
its order one at every prime above the given prime.
The resulting quotient is in $U_E$.
It follows in particular that two ideles have the same
class in $D_T/U_O$ whenever all their differences
$\operatorname{ord}_P-\operatorname{ord}_{\bar P}$
are equal.

For a full self-dual lattice $M\subset E^n$, let $I_M$
be its determinant fractional ideal and write
$d_P(M)=\operatorname{ord}_P(I_M)$.
These orders have finite support. The paired localization
of Hermitian duality, proved in NB1, gives
\[
                 d_{\bar P}(M)=-d_P(M).              \tag{HF.33}
\]
Explicitly, a local basis matrix at $\bar P$ has dual
basis $(\bar B^t)^{-1}$ at $P$; self-duality identifies
this dual module with $M_P$. Its determinant has
the negative conjugate order. This uses local bases,
not a global freeness assumption.

Put $\Delta_P=d_P(M')-d_P(M)$. The source representative
for the relative determinant is
\[
 a(M,M')=\prod_P
       \operatorname{loc}_P(\pi_P)^{\max(\Delta_P,0)}.
 \tag{HF.34}
\]
The nonnegative exponent is deliberate: it is the
source's integer \texttt{toNat}. By (HF.33),
$\Delta_{\bar P}=-\Delta_P$, so
\[
 \operatorname{ord}_P a(M,M')
       -\operatorname{ord}_{\bar P}a(M,M')=\Delta_P.
 \tag{HF.35}
\]
At a fixed prime (HF.33) forces $\Delta_P=0$.
For three lattices the differences $\Delta_P$
telescope. The characterization of $K_DU_E$ therefore
proves
\[
 \operatorname{relDet}(M,M)=1,\qquad
 \operatorname{relDet}(M,M')
   \operatorname{relDet}(M',M'')
                       =\operatorname{relDet}(M,M'').
 \tag{HF.36}
\]
If $M'=\gamma M$ for a rational unitary matrix, then
$\Delta_P=\operatorname{ord}_P(\det\gamma)$.
Choose $k$ by (HF.31) for $z=\det\gamma$.
The equality $z=k/\bar k$ gives
$\Delta_P=\operatorname{ord}_P(k)
-\operatorname{ord}_{\bar P}(k)$, so (HF.35) proves
\[
            \operatorname{relDet}(M,\gamma M)
                 =[\,e_1(\det\gamma)\,]\in D_T/U_O.
 \tag{HF.37}
\]
Rank zero causes no exception: the determinant ideal
is the unit ideal and every displayed order is zero.

Take $\mu_{\mathrm{til},j+1}$ to be the descended
finite character $\mu_{0,j}$ in the source's range,
and one outside it. Together with $D_T,U_O,\chi_D,e_1$
and (HF.34), HF14 and (HF.36)--(HF.37) give every
field of the actual \texttt{trueDtPack}. Its relative
determinant is the source's positive-part representative,
not a replacement by a signed product with different
conventions.

The columns use the finer quotient obtained by imposing
principal-unit congruences also at a prescribed finite
set $Q$ and its conjugates. Call its subgroup $K_{D,Q}$.
The proofs in HF14 apply verbatim with this larger finite
congruence set: it is open, contains the infinite ideles
and the diagonal of $F^\times$, and the same finite
characters kill it. Moreover
$K_{D,Q}U_E=K_DU_E$, since both are the balanced-order
subgroup just computed. Thus
$D_{T,Q}/U_{O,Q}\simeq D_T/U_O$, with the identity on
idele representatives. The positive-part formula
(HF.34), (HF.36), and (HF.37) therefore give exactly
\texttt{trueDtPackQ} for every finite $Q$.
This is the determinant package actually inserted in
the column record.

\paragraph{HF16: a common field for the values of this actual frame.}
Apply TL14 to the frame just constructed. For
$\mu_{0,j}$ the infinity exponents are even, so the
normalization parameter is $r=0$; for $\eta$ it is
$r=1$; for $\chi_\Psi$ it is $r=0$. The two finite-order
characters and the finite residue characters require
only cyclotomic extensions. The finite compositum and
its normal closure therefore give a number field
$M\subset\mathbb C$, Galois over $\mathbb Q$, containing
all host embeddings, all $\omega_k$ values, and
\[
 \begin{array}{c|ccccc}
 \text{character}&\mu_{0,j}&\eta&\chi_\Psi&\delta&\delta'\\ \hline
 \text{normalization }s&0&1/2&0&0&0
 \end{array}
 \qquad
       \chi(P)\,N(P)^{-s}\in M
                    \quad\text{whenever }\chi_P\text{ is unramified}.
 \tag{HF.38}
\]
Here $N(P)^{-s}$ is the positive real normalization
specified in \texttt{valTw}. The assertion is uniform
over all unramified primes; it is not only separate
algebraicity for each prime. TL14 proves the needed
uniformity using finitely many ideal-class representatives
and the parity of the infinity exponents.

Choose a primitive generator $\alpha$ of $M$.
If an actual coefficient embedding $\sigma:L_0\to\mathbb C$
contains $\alpha$ in its image, then it contains $M$:
all elements of $M=\mathbb Q(\alpha)$ are rational
functions of $\alpha$ with nonzero denominators.
Equation (HF.38) then gives the literal
$\texttt{ValLL0}$ when $L_0$ is the fixed coefficient
subfield, and \texttt{ValLL} by its inclusion into the
full coefficient field. For every later assembly the
frame \texttt{Fr.at} uses these same characters and
weights, so the same witness proves the uniform
assertion over those assemblies. Containment of
$\alpha$ is a field choice or a field-containment
hypothesis; it does not follow from the record name
\texttt{HFrame}.

\paragraph{HF17: extending coefficients without changing the preliminary datum.}
We give the finite construction behind
\texttt{BaseExt.exists\_base}, since the original
existence join needs an actual base, not just a possible
field of character values. Fix a preliminary datum
$P_0$ and three finite sets of algebraic complex numbers
$S_0,S_1,S_2$ to be included respectively in $K_s,L_0,L$.
Adjoin their elements, all embeddings of a rational
basis of $E$, the finitely many required roots of
unity, and the finite principalization sets described
below. Take the splitting field over $\mathbb Q$
of the product of their minimal polynomials and
$X^2+1$. This is one finite Galois field $L\subset\mathbb C$,
closed under conjugation and containing $i$.

For ordinary ideal capitulation, choose one integral
ideal $\mathfrak c$ in each ideal class of $E$ and
$\mathfrak c^h=(\gamma_\mathfrak c)$, where $h>0$
is the class number. Adjoin a root
$b_\mathfrak c^h=\phi_0(\gamma_\mathfrak c)$.
Comparison of prime-ideal exponents after extension
shows $\phi_0(\mathfrak c)\mathcal O_L=(b_\mathfrak c)$.
Every other ideal in this class differs by a principal
fractional ideal, so these finitely many roots
principalize every extended ideal. The zero ideal
has generator zero.

The source also needs type-norm generators with a
specified absolute value at \emph{every} embedding of
$L$. For a nonnegative exponent function $r_\tau$ on
all complex embeddings of $E$, suppose
$r_\tau+r_{\bar\tau}=n$ is constant, and put
\[
 \Phi_r(\mathfrak a)=\prod_\tau
           (\tau\mathfrak a\,\mathcal O_L)^{r_\tau},
 \qquad
 N_r(x)=\prod_\tau\tau(x)^{r_\tau}.
 \tag{HF.39}
\]
Adjoin for every class representative a root
$b_\mathfrak c^h=N_r(\gamma_\mathfrak c)$.
Unique factorization of fractional ideals gives
$\Phi_r(\mathfrak c)=(b_\mathfrak c)$.
For any embedding $\sigma:L\to\mathbb C$, its action
permutes the embeddings of $E$ and commutes with
conjugation on that CM field. Pairing conjugate
embeddings in the norm product gives
\[
        |\sigma N_r(x)|^2=|N_{E/\mathbb Q}(x)|^n.
 \tag{HF.40}
\]
It follows by taking the positive $h$th root that
$|\sigma b_\mathfrak c|^2=N(\mathfrak c)^n$.
If $(x)\mathfrak a=(y)\mathfrak c$, then
\[
 \beta_\mathfrak a=N_r(x)^{-1}N_r(y)b_\mathfrak c,
 \qquad
 (\beta_\mathfrak a)=\Phi_r(\mathfrak a),\qquad
 |\sigma\beta_\mathfrak a|^2=N(\mathfrak a)^n.
 \tag{HF.41}
\]
The last equality follows by taking ideal norms of
the relation and using (HF.40). No assertion that
an arbitrary embedding of $L$ commutes with its
conjugation is needed: the conjugation identity
used in (HF.40) is only on the embedded CM field $E$.

Apply this to the finitely many source exponents
$r_\tau=p_j$ on the CM-type and $q_j$ on its conjugate,
$0\le j\le m$. This supplies the exact
\texttt{kPsi} generators and their norm equations.
The ordinary roots supply \texttt{capit}.
Now set the actual coefficient field to $L$, and
choose $K_s=L_0=L$ as subfields of it, with the
inclusion into $\mathbb C$ and the restricted complex
conjugation. Insert the specified host embedding,
roots, capitulation and type-norm generators.
Every field of \texttt{Base} is thereby supplied.
The preliminary part is \emph{equal} to $P_0$:
none of its host, spaces, words, weights, type
primes, masses or windows is changed.
This proves the literal three-set extension
statement and makes the subsequent finite-field
requirements simultaneous.

\paragraph{HF18: original canonical columns with pins and algebraic jets.}
Start with the sized preliminary datum and its actual
nonempty decoration supplied by IG10. In particular it
has \texttt{HasLetters}, diagonal spaces, the exact
$\operatorname{duStarOf}$ values, and the integral
weight compatibility \texttt{MUQ}. HF13 constructs a
frame $f$ for this host and these weights.
Orient each auxiliary factor by IG11 using the
\emph{chosen} primes $f.P$: replacing a factor by
its conjugate preserves its norm and congruences,
and makes its conjugate vanish exactly above the
chosen quadratic prime. Apply the normalized-space
construction of IG11 and SE15. It retains the host,
weights, type primes and letters, and supplies
\[
 i_X'(e_a)=e_a,\qquad
       c_u\,\pi_{T,u}=\pi_{T,u_0}.                  \tag{HF.42}
\]
The auxiliary orientation remains unchanged.
Write $P_0$ for this now fixed preliminary datum,
$X$ for its decoration, and $f$ for the same frame.
Only after this normalization do we extend the
coefficient field.

Let
\[
 \Xi=\{1,\delta,\delta'\},\qquad
 S=\{p\in\operatorname{offSet}(P_0):
                     p\ne p_\delta,\ p\ne p_{\delta'}\}.
 \tag{HF.43}
\]
The type primes, $g_0,g_1$, and every rational prime
dividing the host conductor belong to $S$.
For the conductor primes this follows directly
from $p_\delta,p_{\delta'}\equiv1$ modulo the
conductor; the other exclusions are part of the
initial prime separation. HF12 shows that all three
characters in $\Xi$ are \texttt{IsTwist} for $S$.
In particular this assertion concerns the source's
finite tame twists, not its stronger prime-power
\texttt{IsPsi} predicate.

Choose the following finite algebraic sets before
forming the base. Include the primitive generator
$\alpha$ of HF16 in the fixed coefficient field.
Include the PA13 generator for \texttt{Per.FieldOK},
square roots of $g_0,g_1$, and the TL15
\texttt{Holds} generators for $\Xi$ with exponent
$2m$. Finally include the actual finite set
\[
 Z=\texttt{TwIdxEx.gens}(P_0,f.\chi_\Psi,\Xi)
                                                               \tag{HF.44}
\]
from CA12. Its definition depends on this preliminary
datum, frame character and finite family, not on the
base about to be chosen. HF17 now supplies one base
$B$ with $B.\mathrm{toPre}=P_0$ containing all these
sets. Thus the value-field, roots, period field and
jet-field requirements hold simultaneously. There
is no later enlargement that could change the
column input under discussion.

Choose a prime selector over $F$, and use the
specific ambient datum $J=\texttt{amb0}(B.\mathrm{host})$,
whose full laws were proved in PA9. The TL15
character construction gives \texttt{Chars.Ok},
\texttt{Tied} and \texttt{WeilTied} for this $f$.
Here the orientation of $\pi_G$ is essential: it
identifies the source's selected set $v_G$ with
$f.P.\mathrm{Div}$, so HF7's inverse residue powers
are the local characters required by the constructor.
The reference element $\xi_0$ is a unit at these
primes because $g_k\ne d$ and the host different
formula has no denominator there.

For each $\xi\in\Xi$, its finite character values
lie in the chosen coefficient field. The tied
tuple construction of TL15 reads this same complex
character. If a tuple prime $P$ lies in $Q_E$,
the local reading law gives a ramified prime
$P'$ of $\xi$ with $N(P)=N(P')$. The split
degree-one property makes this norm a rational
prime; it is outside $S$, so it is different from
$d$. Hence $\xi_0$ is a unit at $P$ as well.
The square-root and primitive-root requirements
are precisely those already included by
\texttt{Holds}. CI15--CI17 then give the actual
local indices and constants, with
\[
 M=A=B=m,\qquad D(P)=2m
                  \quad(P\notin v_{G_0}\cup v_{G_1}). \tag{HF.45}
\]
Thus every field of the literal raw
\texttt{InEx.exists\_colIn} is supplied: its index
set is $\Xi$, its tuple is the coefficient-field
reading just constructed, its selected characters
are $1,\delta,\delta'$, and its frame, selector and
ambient datum are the fixed $f$, selector and $J$.
The remaining \texttt{Ok} fields are the already
fixed diagonal-space, prime-containment and
$\operatorname{duStarOf}$ identities.

Apply the actual canonical replacement of CI18
to this raw input. The depth hypothesis needed
there is exactly (HF.45). It preserves the base,
decoration, frame, character family, selector,
ambient datum and \texttt{Ok}; it gives
\texttt{IsTw} and changes the two modulus exponents
to $A=B=m+1$. CA12, with the very same set (HF.44),
now proves \texttt{JetsIn}. Its modulus is
\[
 \left(\prod_{Q\in\mathrm{rowSG}}NQ\right)^{m+1}
 \left(T_u\left(\prod_{Q\in\mathrm{rowSG}}NQ\right)^{m+1}\right),
 \tag{HF.46}
\]
not the raw modulus with exponent $m$.
The two equalities (HF.42) give the literal
\texttt{Pins}, independent of this replacement.

We have therefore constructed one original tuple
$(B,X,I)$ with
\[
 \begin{gathered}
 \texttt{HasLetters}(B.\mathrm{toPre}),\quad
 I.\texttt{Ok},\quad\texttt{Pins}(B,X),\quad
 I.\texttt{IsTw},\quad I.\texttt{JetsIn},\\
 I.C=f,\qquad I.S=S,\qquad
 \operatorname{range}(I.\xi)=\Xi,\qquad
 I.J=\texttt{amb0}(B.\mathrm{host}).                  \tag{HF.47}
 \end{gathered}
\]
HF16 also gives \texttt{ValLL0} for every assembly
over this base, and PA13 gives its actual
\texttt{Per.FieldOK}. Finally $\mu_\star\ne\mu_0$:
otherwise the reading law would give $\delta=1$;
at any prime above $p_\delta$, HF12 with exponent
one would then imply $15\mid1$. The same argument
works for $\delta'$, and the disjoint ramification
distinguishes the two nontrivial twists.

\paragraph{HF19: laws, later stages, twisted forms and all lattice jets on that base.}
Apply TL13 to the \texttt{Ok} input in (HF.47).
This constructs its actual \texttt{Laws}; there
is no additional law-existence premise. The
step-zero class bijection preserves the entire
decorated lattice, and the level-zero stabilizers
are trivial. Consequently a coefficient-field
form at step zero is simply a function on the
finite class set with values in the common
polynomial representation. Pullback along this
class bijection is a linear isomorphism. It
preserves the integral lattice because the
lattice at each matched class is equal; its
complex embedding is exactly the source's
\texttt{e0Tw}. This is the \texttt{Law0} proved
in TX10, with the present input and every twist.

The determinant package needed by
$C=I.\mathrm{cols}(l)$ is supplied by HF15 for the
actual finite set $I.SQ$. Its determinant
characters and lattice-determinant representatives
are consequently the ones used in this record.
The class sums give its Hermitian positive
metrics, and the terminal class spaces have
finite dimension, as in FC7 and SE7--SE8.
LS22 can therefore be applied with precisely
this $B,C$ and its retained \texttt{HasLetters}.
It constructs
\[
 Z=(W,V,S_{\mathrm{st}},N,G):
                  \texttt{Later}(B,C),\qquad
 D=I.\mathrm{rec7}(l,Z).
 \tag{HF.48}
\]
All five stage types in this statement are the
actual types over $B,C$; no pre-existing unrelated
stage tuple is substituted.

For these stages define the twisted forms by the
actual chain and metric forms for each $I.\iota$.
Use the preceding class-pullback isomorphism as
the degree-zero map, and the identity as the
map at the reference tuple $\mu_0$. The lattice
and embedding clauses are exactly the two
statements just proved for \texttt{Law0}, or
identities at $\mu_0$. This fills every field of
$I.\texttt{tw}(l,\mathrm{Law0},Z)$ and proves
\[
           \operatorname{Nonempty}(\texttt{TwForms}(D)).
 \tag{HF.49}
\]
In particular the existential
\texttt{TrueColsTest.exists\_all} is now realized
with the additional canonicality and jet clauses
in (HF.47).

The same input supplies the ordinary kernel
\texttt{Hyp} over the canonical seam of
\emph{any} stages above $B$. Indeed its definite
field and prime data are the \texttt{Ok} data;
the ambient curve and semicharacter are PA9 with
the actual $\chi_\Psi$ of HF8; the character
values follow from HF16; the selector and
forbidden-prime inclusions are those of HF18;
and the local constants and their $\mu_q$ tie
are the CI15--CI17 constants. These are all the
clauses of \texttt{hyp\_of\_colIn}.

Here is the precise passage from \texttt{JetsIn}
to the source's \texttt{JetsAll} for self-dual
lattices. Put $H=\operatorname{homAt}(S,N,N')$.
A coefficient function of the required range
is supported on $c_A^{-1}H$ and invariant under
$c_BH$. If it is nonzero at $x$, then $c_Ax\in H$,
and hence $(c_Ac_B)x\in H$. The global jet
statement in (HF.47) applies to this $x$ with
translation denominator $c_B$; $c_B>0$ follows
from the positive type prime and positive local
norms in (HF.46). The multi-index degree it tests
is
\[
 \deg_\tau=\sum_k(\nu_{u,i+1,\tau})_k,
 \tag{HF.50}
\]
which is exactly the required degree by TL19.
Conversely the indicator of $c_A^{-1}H$ is
invariant under $c_BH$, since $c_BH\subset
c_A^{-1}H$, so testing all such coefficient
functions recovers each prescribed point jet.
This proves the literal range/degree equivalence
and then \texttt{jetsAll\_of\_colIn}, for every
$u,a$ and $i<m$. It supplies ordinary lattice
jets over the same ambient datum and seam;
it does not assert the global trace or spectral
identities used by later theta-test theorems.

\paragraph{HF20: the exact remaining global-core join.}
The conclusions (HF.47)--(HF.49) remove the former
existence boundary at the actual host frame and
at the original \texttt{Ok}/\texttt{Pins}/jets/laws/
\texttt{Later} join. They do not supply all fields
of \texttt{V4.Core}. That record refers to one
specified context $K$, one stage tuple $X$,
one frame $f$, and one nontrivial twisting datum
$T$ on their common base. Its additional
requirements include the common object
\texttt{Obj2}, the column and geometric lattice
layouts, the actual row layout and metric
stability, the theta input/output and gauge
identities, the real and signature ties, the
level-character generators, and the arithmetic
choice ties at the prescribed primes.
None can be inferred solely from nonemptiness
of the component record types.

For clarity, the prime-forcing field already
gives a concrete obstruction to using the
untwisted empty model. Its prescribed finite
set can be chosen explicitly as
\[
 \{g_0,g_1,d\}\ \cup\
       \{0,1,\ldots,\lfloor X''_K(n)\rfloor_{\ge0}\}.
 \tag{HF.51}
\]
Every prime at most $X''_K(n)$, and every
distinguished prime required by \texttt{forced},
lies in this set. But for
$T=\texttt{TwDat.untw}(B,\varnothing)$ the set
$T.\mathrm{pl}$ is empty. The \texttt{forced}
clause of a proposed core would put the prime
$g_0$ into that empty set, a contradiction.
Thus the original three-character column
construction, or LS22's bare empty local/twist
stage, cannot by itself be promoted to a core
for that empty twisting datum.

The exact constructive reduction is instead
the source's \texttt{coreOfLay}. On one fixed
$X,f,T$, retain $K.\texttt{Over}(X.B.\mathrm{toPre})$
and supply its actual layout $Y$,
\texttt{RecOK}, the universally quantified
\texttt{LedOK}, \texttt{MuIsW}, \texttt{forced},
the numerical gate identity, the untwisted
row-kernel identity, metric stability, real
tie, true signature, primitive level-character
generators, and the specified arithmetic-choice
conductor-support and exponent equalities
$Y.\mathrm{nch}.Q(w)=T.Q(w)$ and
$Y.\mathrm{nch}.k'(w)=T.k(w)$. Then its common object
is $\operatorname{layObj}(Y)$; the seven column
fields come from $Y$ and \texttt{RecOK}; the
geometric filling and ledger fields come from
$Y$ and \texttt{LedOK}; its row layout is the
given untwisted row identity; and its character,
gauge and arithmetic ties are the stated
same-object hypotheses. In particular the
off-twist value clause uses the \emph{same}
\texttt{MuIsW.off} equality inside the
\texttt{LedOK} conclusion, rather than a
newly chosen tuple map. This fills the literal
core when those remaining suppliers are
available. HF20 records this exact conditional
join and the empty-model obstruction; it makes
no unconditional core, quantitative window,
theta-test, RH or GRH assertion.
